\section{Related work}
\label{sec:related}

\subsection{Access control for constrained devices}

Role-based access control as formalised by Sandhu et al.\ \citep{sandhu1996rolebased} and standardised by NIST \citep{ferraiolo2001proposed} is still the reference model for multi-stakeholder authorization, and its IoT adaptations share a structural assumption that becomes awkward in the field: some trusted decision component must be reachable when access is requested. RBAC extensions for IoT \citep{ye2014efficient}, attribute-based evaluation \citep{hu2014guide,servos2017current,hernandez2013distributed} and capability-token delegation \citep{hernandez2013distributed} place that trust in a policy engine, an attribute authority and a credential issuer respectively, retaining a central enforcement point \citep{ye2014efficient,servos2017current}, distributing revocation into tokens rather than solving it \citep{hernandez2013distributed}, or splitting the model between humans and devices while leaving the role store centralised \citep{alshehri2016access}. Sciancalepore et al.\ \citep{sciancalepore2017oauthiot} quantify what open-standard authorization stacks cost on constrained hardware, the empirical case for moving policy evaluation off the sensor entirely; Novo \citep{novo2018blockchain} argued early that the role store itself belongs on a distributed ledger. We take that as given and ask what it costs.

\subsection{Blockchain-based access control and how it is evaluated}

Capability-token systems on public chains established feasibility without practicality: FairAccess \citep{ouaddah2016fairaccess} encodes grants as Bitcoin-like transfers with confirmation windows measured in minutes, and BlendCAC \citep{xu2018blendcac} uses ERC-721 tokens with per-transaction fees no sensor duty cycle can absorb. Neither addresses role hierarchy, geographic scoping or seasonal expiry, which agriculture actually needs.

Permissioned platforms remove the latency and fee obstacles, and Fabric-based IoT work has grown accordingly: testbed architectures and performance models \citep{shih2022blockchainbased,pajooh2021hyperledger,ke2023performance}, trust and reputation layers \citep{putra2021trust}, contract-enforced domain policy on cloud-hosted nodes where memory and thermal limits do not apply \citep{usman2024blockchainbased}, and a data marketplace at scale \citep{gonzalez2023use}. All report aggregate throughput and latency from a controlled setting. Two are closer to our method: Idrissi and Palmieri \citep{idrissi2024blockchain} give an end-to-end Fabric latency breakdown on constrained ARM hardware in a healthcare context, the nearest methodological precedent for the partition in Section~\ref{sec:results-latency-accounting}, and Tzenetopoulos et al.\ \citep{tzenetopoulos2022hlfkubed} establish a Raspberry Pi 3B+ reference point discussed in Section~\ref{sec:discussion}.

The supplementary material's system-comparison table summarises how these and further systems \citep{zaidi2023fabrication,bandara2021fabric} are evaluated. The pattern is consistent: laboratory testbeds or simulation, runs of hours or days, aggregate reporting. We have not found prior work that times an authorization-associated span separately from the rest of the write path on a field deployment of this duration, or that reports revocation behaviour over a multi-week window; that combination, rather than any single mechanism in Section~\ref{sec:design}, is the contribution.

The methodological gap is not merely duration. Even where a Fabric-based IoT study reports latency broken down by pipeline stage \citep{idrissi2024blockchain} or peer count \citep{ke2023performance,pajooh2021hyperledger}, the reported figures come from a single controlled run or a small number of repetitions without an explicit design against order confounds; none of the surveyed systems states, and then tests, whether configuration and execution sequence were deliberately separated. Section~\ref{sec:methodology} treats that separation as a first-class methodological requirement rather than an incidental design choice, because the alternative -- an ascending or otherwise monotonic execution order confounded with the factor of interest -- is precisely the pattern that produced the single-run, single-block campaigns this article's own earlier internal drafts had to redesign (Section~\ref{sec:stats-treatment}). Formal-methods work adjacent to this literature \citep{lamport2002specifying,newcombe2015how} checks a specification against itself, not an implementation against its specification; the audit in Section~\ref{sec:audit} is closer in spirit to the latter, and we are not aware of a published Fabric-based access-control study that pairs a multi-week field trace with a code-to-specification audit of the kind reported here.

\subsection{Agricultural IoT and edge infrastructure}

Blockchain in agriculture has concentrated on traceability after harvest \citep{tian2017supply,lin2019food,bandara2021fabric}, where the unit of interest is a shipment rather than a zone-scoped actuation. Safeer et al.\ \citep{safeer2025agrifarming} combine computer vision, IoT and blockchain for climate-smart cultivation and, like most pre-harvest work, treat authorization as an implementation detail rather than a measured component. Liu et al.\ \citep{liu2019distributed} and Wang et al.\ \citep{wang2020convergence} reduce processing latency at the agricultural edge without addressing the authorization layer above it, and Mohanta et al.\ \citep{mohanta2025smartcontract} apply smart-contract enforcement in smart grids, a neighbouring domain with comparable device constraints and a similarly centralised starting point.

On identity infrastructure, W3C Decentralized Identifiers \citep{w3c2022did} and gateway-assisted verification for constrained devices \citep{ferdous2019search,tara2023decentralized} offer an eventual migration away from X.509, treated here as future work rather than a deployed alternative. On the radio side, Adelantado et al.\ \citep{adelantado2017understanding} characterise the LoRaWAN capacity limits that bound achievable sensor density per gateway, a constraint visible in the energy results (Section~\ref{sec:energy-results}).
